\subsection{Literal braid words for the modular CM pullback}

\begin{result}{A.4.4}{Proved}{An explicit modular embedding with its peripheral orientation}
Write
\[
 L=B_4/\langle z_4\rangle,\quad
 q:L\longrightarrow B_3/\langle z_3^2\rangle,
 \quad q(a)=q(c)=a,\quad q(b)=b,
\]
where $(a,b,c)=(\sigma _1,\sigma _2,\sigma _3)$,
$z_3=(ab)^3$ and $z_4=(abc)^4$.  Put
$u=abc$ and $d=a^2bc$.  There is an explicitly marked isomorphism
\[
 \rho_+:B=\langle s,r\mid s^2=r^3=1\rangle
       \xrightarrow{\ \sim\ }q^{-1}(\langle ab\rangle),
 \qquad s\longmapsto u^2,
 \quad r\longmapsto d^{-1},
 \quad \tau=sr\longmapsto abc a^{-1}.             \tag{FCE.1}
\]
On the cyclic quotient $C_6=\langle ab\rangle$ it sends $\tau$ to
$ab$.  Its full profinite completion is the isomorphism onto the
actual preimage of this finite subgroup in $\widehat L$.

For the Ferrari basis
\[
 x=ca^{-1},\qquad y=aba c^{-1}b^{-1}a^{-1}
\]
of $K=\ker q$, its restriction to $B'$ is given by
\[
 p=s\tau^{-3}\longmapsto xy,
 \qquad p_1=\tau p\tau^{-1}\longmapsto y^{-1}xy.
                                                               \tag{FCE.2}
\]
Its peripheral orientation and based conjugator are fixed by
\[
 \rho_+(\tau)^6=(xy)^{-1}[x,y]^{-1}(xy)
       \quad\hbox{in }L,\qquad [x,y]=xyx^{-1}y^{-1}.             \tag{FCE.3}
\]
\end{result}
\begin{proof}
The actual braid identities $u^4=d^3=z_4$ establish the two torsion
relations in (FCE.1).  Direct cancellation gives
$u^2d^{-1}=abc a^{-1}$.  Folding sends $u^2$ to $z_3=(ab)^3$
and $d$ to $a^2ba=(ab)^2$, so the induced map to the order-six
subgroup is exactly the abelianization $B\twoheadrightarrow C_6$,
with $\tau\mapsto ab$.

The free-kernel presentation in E.261 gives
\[
 \operatorname{Ad}(a):(x,y)\longmapsto(x,yx),\qquad
 \operatorname{Ad}(b):(x,y)\longmapsto(xy^{-1},y).
\]
It follows directly, or by substitution in the literal braid words,
that
\[
 \rho_+(\tau)=y^{-1}(ab),\qquad
 \rho_+(s)=x^{-1}y^{-1}(ab)^3,
\]
and then the two identities in (FCE.2) follow.  Their right sides
are a free basis of $K$: if they are denoted by $v,w$, the inverse
substitutions are $y=vw^{-1}$ and $x=vwv^{-1}$.

The index-six subgroup $B'$ is torsion-free and free of rank two,
as follows from its graph-of-groups cover of $C_2*C_3$ (Euler
characteristic $6(1/2+1/3-1)=-1$).  Thus the surjection
$B'\to K$ just obtained is an isomorphism: after arbitrary free-basis
identifications it is a surjective endomorphism of a finitely
generated free group, which is Hopfian.  The map in (FCE.1) is now
an isomorphism on both the kernels and the quotient $C_6$, proving
its assertion.

For completeness the exact peripheral word can be computed before
taking the central quotient.  Ferrari combing of
$\rho_+(\tau)^6z_4^{-1}$ gives $y^{-1}x^{-1}yx$, with trivial
three-strand factor.  Free reduction identifies that word with the
right side of (FCE.3).  The program
\texttt{check\_quartic\_embedding.py} independently verifies every
braid identity above in the faithful Artin action of $B_4$ on $F_4$;
it also performs the Ferrari combing and checks the full cyclic
commutator orientation.

\par\medskip\noindent\textbf{The actual geometric pullback and its homology marking.}\quad
This also identifies the image of the geometric division-quartic
map, without using E.260 as an additional assertion.  Put
$D=4a_0^3+27b_0^2$ (the subscripts distinguish these coefficients
from the braid letters).  Rigidification of the original elliptic
involution gives weights $(2,3)$ on $(a_0,b_0)$.  Since $D$ has
weight six, the slice $D=108$ gives an actual isomorphism of stacks
\[
 [\{D\ne0\}/\mathbb G_m{}_{(2,3)}]
   \simeq[\{D=108\}/\mu_6]
   \xrightarrow{\ \sim\ }
 [ (\mathcal E_0\setminus\{O\})/\mu_6],
 \quad \mathcal E_0:Y^2=Z^3+1,                         \tag{FCE.4}
\]
where the second map and its inverse are the literal formulas
\[
 (Z,Y)=(-a_0/3,-b_0/2),\qquad (a_0,b_0)=(-3Z,-2Y).
                                                               \tag{FCE.5}
\]
Indeed $D=108(Y^2-Z^3)$, and the residual action is
$(Z,Y)\mapsto(v^2Z,v^3Y)$ for $v\in\mu_6$.
The first map is the scalar-slice quotient: the six choices of a
sixth root normalizing $D$ differ precisely by this displayed
$\mu_6$ action.  Thus it is a stack isomorphism, not merely a map
of coarse spaces.  The source is the orbifold affine line with
an order-three point at $a_0=0$, an order-two point at $b_0=0$,
and the discriminant cusp removed.  Its orbifold fundamental
group is the displayed $C_2*C_3$.

The point in (FCE.5) is the ordered second point in the monic
quartic construction: the positive pole is $O$ and the negative
pole is $(-a_0/3,-b_0/2)$.  This follows directly from
$Z=(t+x^2+a_0)/2-a_0/3$ and
$Y=((t+x^2+a_0)/2)x+b_0/2$ in the mutually inverse quartic formulas.
Consequently (FCE.4) is the actual pullback of the ordinary
universal curve $\mathcal M_{1,2}\to\mathcal M_{1,1}$ along
$B\operatorname{Aut}(\mathcal E_0,O)=B\mu_6$.
Its fundamental-group map is an isomorphism on the fibre
$\pi_1(\mathcal E_0\setminus\{O\})=F_2$ and has the indicated
$\mu_6$ quotient.  The ordinary Birman sequence therefore
identifies its image with the full inverse image of this very
CM subgroup in $\Gamma_{1,2}$.

Here is the homology marking of that subgroup.  Set
$\zeta=e^{2\pi i/3}$ and $v=e^{\pi i/3}=1+\zeta$.
Use the oriented period basis $(1,\zeta)$ of the hexagonal
uniformization of $\mathcal E_0$, identified with the oriented
pointpush basis $(\xi_1,\xi_2)$ of (E.261).  This is an actual
marking of the torus.  The displayed $v$-action multiplies the
invariant differential $dZ/Y$ by $v^{-1}$ and hence acts on its
uniformizing coordinate by $v^{-1}=-\zeta$.  On this basis its
matrix is
\[
 1\longmapsto-\zeta,\qquad
 \zeta\longmapsto1+\zeta,\qquad
 \begin{pmatrix}0&1\\-1&1\end{pmatrix}
 =UL_0=Q(ab).                                           \tag{FCE.6}
\]
The same marking is available in the Teichm\"uller space of the
once-punctured torus; transport along it identifies ordinary
forgetting with the folding map $q$.  Thus the CM subgroup is
literally $\langle ab\rangle$ in
$B_3/\langle z_3^2\rangle$, not an abstractly isomorphic cyclic
subgroup.  Formula (FCE.4) and the kernel--quotient comparison
identify the geometric image with $q^{-1}(\langle ab\rangle)$.
Any remaining based choice in its fibre is the actual inner
element retained in A.4.5; no assertion that this element is
trivial is needed.

Finally the completion of the discrete pullback has kernel
$\widehat B'\cong\widehat F_2$, since $B'$ has finite index in $B$.
The full Birman injection identifies this kernel with the actual
closed kernel $\widehat K$ of $\widehat L\to\widehat{B_3/\langle
z_3^2\rangle}$.  The cyclic subgroup survives in the latter group,
for example in its integral homology representation modulo three
and four.  Therefore the completed map has this same kernel and
quotient, and is an isomorphism onto the actual pullback.
\end{proof}

This gives specified finite words for pulling a finite quotient of
$L$ back to $B$.  In particular the quotient in A.4.3 is invariant
under $L$, because both its parity map and its signed matrix map
are equivariant under braid conjugation; its normal kernel therefore
defines a finite quotient of $L$ of order
$24\cdot768\cdot3^{114}$, to which (FCE.1) applies literally.
Comparison with the geometric division-quartic map of E.260 uses
its actual peripheral orientation and basepoint transport; (FCE.3)
records those data for $\rho_+$ rather than assigning an unspecified
isomorphism the same marking.  No finite parity exclusion follows
merely from the existence of (FCE.1).
